% !TeX program = xelatex
\documentclass[aspectratio=169]{beamer}

\usepackage{import}

\subimport{../../}{system.tex}
\UseTemplateSet{
  layout = beamer,
  globalfonts = {cmu,shanggu},
  fonts = {noto,phoenician,coptic,armenian,georgian,anatolian,old_italic,runic,glagolitic},
  features = {tables,image}
}

\newcommand{\NotoLatinSans}[1]{{\textsf{\NOT{#1}}}}

\newcommand{\GK}[1]{{\NotoLatinSans{#1}}}

\title{希臘系文字}
\subtitle{世界文字總覽~第一講}
\author{Kelvin Yang}
\date{March 7, 2026}

\begin{document}
\begin{frame}[plain]
    \titlepage
\end{frame}
\AgendaFrame

\section{希臘系文字}
\begin{frame}{希臘系文字}
    本講對希臘系文字進行概述，並重點介紹希臘文、拉丁文以及西里爾文。
    \begin{block}{希臘系文字}
        \textbf{希臘系文字}（Hellenic scripts），即包含希臘文以及相關衍生文字在內的、位於希臘--羅馬文化圈影響範圍內的諸多族羣所使用的文字系統。
    \end{block}

    總體上，希臘系文字可以分爲三個子系：
    \begin{itemize}
        \item \textbf{義大利子系}（Italic）
        \item \textbf{斯拉夫子系}（Slavic）
        \item \textbf{安納托利亞子系}（Anatolic）
    \end{itemize}
    除了這三子系之外，還有一些相對零散、獨立，如科普特文、亞美尼亞文、格魯吉亞文、高加索阿爾巴尼亞文等。
\end{frame}

\begin{frame}{希臘早期歷史}
    希臘的早期歷史大概可以分爲以下階段：
    \begin{itemize}
        \item \textbf{邁錫尼時期}（Mycenaean Greece）：約公元前十七--十二世紀，希臘進入邁錫尼文明的形成與發展期。其核心特徵是以要塞化宮殿爲中心的青銅時代晚期社會，並出現線形文字 B 所記錄的最早希臘語。
        \item \textbf{黑暗時代}（Greek Dark Ages）：約公元前十二--九世紀前後，邁錫尼宮殿體系在多因素壓力下崩潰，書寫中斷、長距貿易萎縮，人口與聚落規模整體收縮，社會趨於地方化重組。
        \item \textbf{古風時代}（Archaic Greece）：約公元前八--六世紀，經濟與交換網絡恢復並擴張，城邦與公民共同體逐步成形，同時殖民活動與成文法、公共祭祀等制度化過程推動「泛希臘世界」的定型。
    \end{itemize}

    希臘文並不是希臘語的第一種文字。早在邁錫尼時期，邁錫尼人就使用\textbf{線形文字B}（Linear B）來書寫邁錫尼希臘語。
    線形文字B與後來的希臘文無關，可能來源於線形文字A（Linear A），在邁錫尼希臘衰落後就不再使用了。
\end{frame}

\begin{frame}{希臘文的出現與發展}
    希臘文來源於腓尼基文（Phoenician），大約產生於古風時代初期。目前發現最早的希臘文是公元前八世紀的狄皮隆銘文（Dipylon inscription）。

    \begin{OneImage}[htbp][\linewidth][0.6\textheight]{_assets/20260305_212126.png}[狄皮隆銘文]\end{OneImage}
\end{frame}

\begin{frame}{希臘文的出現與發展}
    早期的希臘文並沒有統一，各地都有自己的標準。公元前403年，伯羅奔尼撒戰爭後雅典決議改用愛奧尼亞字母（Ionic）。
    公元前四世紀，隨着雅典的文化影響力，愛奧尼亞字母也推廣到了整個希臘地區，成爲了標準的希臘字母。
    到公元400年左右，希臘字母才演變爲今天的樣子。
    \begin{block}{從腓尼基文到希臘文}
        希臘人在借用腓尼基文時，對其進行了下列改造：
        \begin{itemize}
            \item 將原本用來表示喉音（pharyngeal）的\PH{𐤀}、\PH{𐤄}、\PH{𐤇}、\PH{𐤏}和表示近音（approximant）的\PH{𐤅}、\PH{𐤉}六個字母用來表示元音，即\NotoLatinSans{Α}、\NotoLatinSans{Ε}、\NotoLatinSans{Η}、\NotoLatinSans{Ο}和\NotoLatinSans{Υ}、\NotoLatinSans{Ι}。
            \item 將\PH{𐤈}改爲\NotoLatinSans{Θ}，\PH{𐤎}改爲\NotoLatinSans{Ξ}，\PH{𐤔}表示\NotoLatinSans{Σ}。
            \item 新造字母\NotoLatinSans{Φ}、\NotoLatinSans{Χ}、\NotoLatinSans{Ψ}、\NotoLatinSans{Ω}。
            \item 原本將\PH{𐤅}改造爲\NotoLatinSans{Ϝ}，將\PH{𐤑}改造爲\NotoLatinSans{Ϻ}，將\PH{𐤒}改造爲\NotoLatinSans{Ϙ}，最後棄用。
            \item 書寫方向由原來的從右往左變爲從左往右。
        \end{itemize}
    \end{block}
\end{frame}

\begin{frame}{希臘文}
    我們選取公元前五世紀的阿提卡希臘語（Attic Greek）對希臘文進行講解。
    
    \small
    {\setlength{\columnsep}{0.6em}% ← 三欄之間距離：0.4em~0.8em 自己微調
    \setlength{\tabcolsep}{3pt}%  ← 表內欄距：2pt~4pt 自己微調

    \begin{columns}[T]
        \column{0.3\textwidth}
            \centering
            \begin{tabularx}{\linewidth}{@{}lllc@{}}
                \toprule
                \NiceHead{ 字母 & 名稱 & IPA & 來源 \\ }
                \GK{A} \GK{α} & \GK{ἄλφα}   & /a(ː)/ & \PH{𐤀} \\
                \GK{B} \GK{β} & \GK{βῆτα}   & /b/ & \PH{𐤁} \\
                \GK{Γ} \GK{γ} & \GK{γάμμα}  & /g,(ŋ)/ & \PH{𐤂} \\
                \GK{Δ} \GK{δ} & \GK{δέλτα}  & /d/ & \PH{𐤃} \\
                \GK{Ε} \GK{ε} & \GK{εἶ}     & /e/ & \PH{𐤄} \\
                \GK{Ζ} \GK{ζ} & \GK{ζῆτα}   & /dz/ & \PH{𐤆} \\
                \GK{Η} \GK{η} & \GK{ἦτα}    & /ɛː/ & \PH{𐤇} \\
                \GK{Θ} \GK{θ} & \GK{θῆτα}   & /tʰ/ & \PH{𐤈} \\
                \bottomrule
            \end{tabularx}

        \column{0.3\textwidth}
            \centering
            \begin{tabularx}{\linewidth}{@{}lllc@{}}
                \toprule
                \NiceHead{ 字母 & 名稱 & IPA & 來源 \\ }
                \GK{Ι} \GK{ι} & \GK{ἰῶτα}    & /i(ː)/ & \PH{𐤉} \\
                \GK{Κ} \GK{κ} & \GK{κάππα}   & /k/ & \PH{𐤊} \\
                \GK{Λ} \GK{λ} & \GK{λάμβδα}  & /l/ & \PH{𐤋} \\
                \GK{Μ} \GK{μ} & \GK{μῦ}      & /m/ & \PH{𐤌} \\
                \GK{Ν} \GK{ν} & \GK{νῦ}      & /n/ & \PH{𐤍} \\
                \GK{Ξ} \GK{ξ} & \GK{ξεῖ}     & /ks/ & \PH{𐤎} \\
                \GK{Ο} \GK{ο} & \GK{οὖ}      & /o/ & \PH{𐤏} \\
                \GK{Π} \GK{π} & \GK{πεῖ}     & /p/ & \PH{𐤐} \\
                \bottomrule
            \end{tabularx}

        \column{0.3\textwidth}
            \centering
            \begin{tabularx}{\linewidth}{@{}lllc@{}}
                \toprule
                \NiceHead{ 字母 & 名稱 & IPA & 來源\\ }
                \GK{Ρ} \GK{ρ}   & \GK{ῥῶ}      & /r,r̥/ & \PH{𐤓} \\
                \GK{Σ} \GK{σ/ς} & \GK{σῖγμα}   & /s,(z)/ & \PH{𐤔} \\
                \GK{Τ} \GK{τ}   & \GK{ταῦ}     & /t/ & \PH{𐤕} \\
                \GK{Υ} \GK{υ}   & \GK{ὖ}       & /uː/ & \PH{𐤅} \\
                \GK{Φ} \GK{φ}   & \GK{φεῖ}     & /pʰ/ & \PH{𐤒} \\
                \GK{Χ} \GK{χ}   & \GK{χεῖ}     & /kʰ/ \\
                \GK{Ψ} \GK{ψ}   & \GK{ψεῖ}     & /ps/ \\
                \GK{Ω} \GK{ω}   & \GK{ὦ}       & /ɔː/ \\
                \bottomrule
            \end{tabularx}
    \end{columns}
    }
\end{frame}

\begin{frame}{希臘文文本案例}
    \textbf{馬太福音7:7--11}

    \GK{Αἰτεῖτε, καὶ δοθήσεται ὑμῖν· ζητεῖτε, καὶ εὑρήσετε· κρούετε, καὶ ἀνοιγήσεται ὑμῖν. πᾶς γὰρ ὁ αἰτῶν λαμβάνει καὶ ὁ ζητῶν εὑρίσκει καὶ τῷ κρούοντι ἀνοιγήσεται. Ἢ τίς ἐστιν ἐξ ὑμῶν ἄνθρωπος, ὃν αἰτήσει ὁ υἱὸς αὐτοῦ ἄρτον, μὴ λίθον ἐπιδώσει αὐτῷ; ἢ καὶ ἰχθὺν αἰτήσει, μὴ ὄφιν ἐπιδώσει αὐτῷ; εἰ οὖν ὑμεῖς πονηροὶ ὄντες οἴδατε δόματα ἀγαθὰ διδόναι τοῖς τέκνοις ὑμῶν, πόσῳ μᾶλλον ὁ Πατὴρ ὑμῶν ὁ ἐν τοῖς οὐρανοῖς δώσει ἀγαθὰ τοῖς αἰτοῦσιν αὐτόν.}

    你們祈求、就給你們．尋找、就尋見．叩門、就給你們開門。因爲凡祈求的、就得着．尋找的、就尋見．叩門的、就給他開門。你們中間、誰有兒子求餅、反給他石頭呢．求魚、反給他蛇呢。你們雖然不好、尚且知道拿好東西給兒女、何况你們在天上的父、豈不更把好東西給求他的人麼。
\end{frame}


\begin{frame}{零散的希臘系文字}
    亞歷山大大帝（Alexander the Great）征服埃及後，其繼業者托勒密一世（Ptolemy I Soter）在埃及建立了托勒密王朝（Ptolemaic Kingdom）。
    
    此時埃及聖書文已經發展出了世俗體（Demotic），與托勒密王朝的官方文字希臘文並行。這兩種文字都用於書寫埃及語，最終融合爲\textbf{科普特文}（Coptic）。科普特文也是一種全音位文字。
    
    科普特文大多來自於希臘文，但有下列字母來源於世俗體：\CO{Ϣ}、\CO{Ϥ}、\CO{Ϧ}、\CO{Ϩ}、\CO{Ϫ}、\CO{Ϭ}、\CO{Ϯ}。
    \begin{block}{文本樣例}
        \CO{ⲥⲟⲩⲙⲟⲥⲉ ⲣⲱⲙⲉ ⲛⲓⲙ ⲉⲩϣⲏϣ ⲉ ⲛⲉⲩⲉⲣⲏⲩ ϩⲛ ⲟⲩⲇⲓⲕⲁⲓⲟⲥⲩⲛⲏ.}
    \end{block}
\end{frame}

\begin{frame}{零散的希臘系文字}
    在基督教傳播需要本族語譯經的背景下，公元405年前後，亞美尼亞神學家梅斯羅普・馬什托茨（Mesrop Mashtots）主導創制了\textbf{亞美尼亞文}（Armenian）。
    亞美尼亞文明顯是在希臘文的啓發下創制的，也可能受到了敘利亞文等其他文字的影響。

    \begin{block}{文本樣例}
        \HY{Բոլոր մարդիկ ծնվում են ազատ ու հավասար` իրենց արժանապատվությամբ և իրավունքներով:}
    \end{block}

    \textbf{格魯吉亞文}（Georgian）的來源暫不清楚，但應該也受到了希臘文的影響。其出現不晚於五世紀。

    \begin{block}{文本樣例}
        \KA{ყველა ადამიანი იბადება თავისუფალი და თანასწორი თავისი ღირსებითა და უფლებებით.}
    \end{block}
\end{frame}

\section{安納托利亞子系}
\begin{frame}{安納托利亞子系}
    \textbf{安納托利亞子系}（Anatolic）文字意指在公元前五--四世紀左右使用於安納托利亞地區，用以記錄安納托利亞語以及弗里幾亞語的一些全音位文字。
    這些文字之間不一定有演化關係，但在空間上聚集，並明顯與希臘文有關，雖然具體的源流關係仍然不清楚。
    這些文字中最重要的包括呂底亞文（Lydian）、呂基亞文（Lycian）和卡利亞文（Carian）等，他們和希臘文都有着相當明顯的對應關係，但具體音值上卻又常有出入。

    \begin{OneImage}[htbp][\linewidth][0.4\textheight]{_assets/20260306_154235.png}[卡里亞語銘文\CA{𐊨𐊣𐊠𐊦𐊹𐊸} /\GK{qlaʎiç}/]\end{OneImage}
\end{frame}

\section{義大利子系}
\begin{frame}{義大利子系}
    \textbf{義大利子系}（Italic）指希臘文傳播到古義大利半島後，演變出的以伊特魯里亞文（Etruscan）爲代表的古義大利文及其衍生出的一系列文字。

    \textbf{古義大利文}（Old Italic）衍生自埃維亞希臘文（Euboea Greek），一種希臘文統一前的西型變種。
    公元前八世紀之後，使用埃維亞希臘文的希臘殖民者們來到了義大利半島，這裏的人們也就在埃維亞希臘文的基礎上創製出了新的文字來表達自己的語言。
    其中最重要的，是伊特魯里亞人所使用的伊特魯里亞文。
    
    \begin{OneImage}[htbp][\linewidth][0.4\textheight]{_assets/20260307_020029.png}[古伊特魯里亞文（公元前七--五世紀）]\end{OneImage}
\end{frame}

\begin{frame}{拉丁文}
    一般認爲，\textbf{拉丁文}（Latin）出現於公元前七世紀的義大利，由伊特魯里亞文演變而來。拉丁文採取了伊特魯里亞文26個字母中的21個，我們選取公元前25年--公元200年之間的古典拉丁語對拉丁文進行介紹。

    \small
    {\setlength{\columnsep}{0.6em}% ← 三欄之間距離：0.4em~0.8em 自己微調
    \setlength{\tabcolsep}{3pt}%  ← 表內欄距：2pt~4pt 自己微調

    \begin{columns}[T]
        \column{0.3\textwidth}
            \centering
            \begin{tabularx}{\linewidth}{@{}lllcc@{}}
                \toprule
                \NiceHead{ 字母 & 名稱 & IPA & 來源 & 對照 \\ }
                A a & ā & /a(ː)/ & \OI{𐌀} & \GK{Α α} \\
                B b & bē & /b/ & \OI{𐌁} & \GK{Β β} \\
                C c & cē & /k/ & \OI{𐌂} & \GK{Γ γ} \\
                D d & dē & /d/ & \OI{𐌃} & \GK{Δ δ} \\
                E e & ē & /ɛ(ː)/ & \OI{𐌄} & \GK{Ε ε} \\
                F f & ef & /f/ & \OI{𐌅} & \GK{Ϝ ϝ} \\
                G g & gē & /g/ & \OI{𐌂} & \GK{Γ γ} \\
                \bottomrule
            \end{tabularx}

        \column{0.3\textwidth}
            \centering
            \begin{tabularx}{\linewidth}{@{}lllcc@{}}
                \toprule
                \NiceHead{ 字母 & 名稱 & IPA & 來源 & 對照 \\ }
                H h & hā & /h/ & \OI{𐌇} & \GK{Η η} \\
                I i & ī & /i(ː)/ & \OI{𐌉} & \GK{Ι ι} \\
                K k & kā & /k/ & \OI{𐌊} & \GK{Κ κ} \\
                L l & el & /l/ & \OI{𐌋} & \GK{Λ λ} \\
                M m & em & /m/ & \OI{𐌌} & \GK{Μ μ} \\
                N n & en & /n/ & \OI{𐌍} & \GK{Ν ν} \\
                O o & ō & /ɔ(ː)/ & \OI{𐌏} & \GK{Ο ο} \\
                \bottomrule
            \end{tabularx}

        \column{0.3\textwidth}
            \centering
            \begin{tabularx}{\linewidth}{@{}lllcc@{}}
                \toprule
                \NiceHead{ 字母 & 名稱 & IPA & 來源 & 對照\\ }
                P p & pē & /p/ & \OI{𐌐} & \GK{Π π} \\
                Q q & qū & /kʷ/ & \OI{𐌒} & \GK{Ϙ ϙ} \\
                R r & er & /r/ & \OI{𐌓} & \GK{Ρ ρ} \\
                S s & es & /s/ & \OI{𐌔} & \GK{Σ σ/ς} \\
                T t & tē & /t/ & \OI{𐌕} & \GK{Τ τ} \\
                V u & ū & /u(ː)/ & \OI{𐌖} & \GK{Υ υ} \\
                X x & ix & /ks/ & \OI{𐌗} & \GK{Χ χ} \\
                \bottomrule
            \end{tabularx}
    \end{columns}
    }
\end{frame}


\begin{frame}{拉丁文文本案例}
    \textbf{高盧戰記1:1}

    Gallia est omnis dīvīsa in partēs trēs, quārum ūnam incolunt Belgae, aliam Aquītānī, tertiam quī ipsōrum linguā Celtae, nostrā Gallī appellantur. Hī omnēs linguā, īnstitūtīs, lēgibus inter sē differunt. Gallōs ab Aquītānīs Garumna flūmen, ā Belgīs Mātrona et Sēquana dīvidit.

    全高盧被劃分為三個部分，其中一部分由比利時人居住，另一部分由阿基坦人居住，第三部分由那些用他們自己的語言稱為凱爾特人、我們稱為高盧人的人居住。這些人之間在語言、習俗和法律上各有不同。高盧人與阿基坦人之間由加爾諾河劃分，與比利時人之間則由馬特羅納河和塞庫安納河劃分。
\end{frame}

\begin{frame}{盧恩文}
    \textbf{盧恩文}（Runes）的形成大約在公元一--二世紀間，由日耳曼人創制來拼寫日耳曼語，可能受到了北傳的伊特魯里亞文或拉丁文的影響。

    最早的盧恩文是古弗薩克文（Elder Futhark），後來才分化出北歐的新弗薩克文（Younger Futhark）與盎格魯--撒克遜人（Anglo-Saxon）的弗托克文（Futhorc）。

    盧恩文在基督教化以後逐漸被拉丁字母擠出正式書寫領域，西歐約 8--10 世紀已大致完成，北歐則使用到約12--13 世紀。
    不過，也有一些民間書寫傳統到很晚都沒有斷絕，最典型的例子就是瑞士達拉納（Dalarnas）的達拉卡利亞如尼文（Dalecarlian runes），一直使用到二十世紀八十年代。

    \begin{block}{文本樣例}
        \RU{ᛠᛚᛚ ᚠᚩᛚᚳ ᚹᛖᚩᚱᚦᚪᚦ ᚠᚱᛖᚩ ᚪᚾᛞ ᛖᚠᚾᛖ ᛒᛖ ᚪᚱᛖ ᚪᚾᛞ ᚱᛁᚻᛏᚢᛗ ᚷᛖᛒᚩᚱᛖᚾ.}
    \end{block}
\end{frame}

\section{斯拉夫子系}
\begin{frame}{斯拉夫子系}
    \textbf{斯拉夫子系}（Slavic）意指希臘文東傳後，由斯拉夫人改造、使用的一系列拼寫斯拉夫語的文字。
    
    一般認爲，在大約九世紀，拜占庭帝國的傳教士基里爾（Cyril）和美多德（Methodius）出於傳教的需要，創制了\textbf{格拉哥里文}（Glagolitic）。
    格拉哥里文在形態上與希臘文相當不一致，但其底層邏輯和結構基礎與希臘文相當一致。

    \begin{block}{文本樣例}
        \GL{ⰲⱐⱄⰻ ⰱⱁ ⰾⱓⰴⰻⰵ ⱃⱁⰴⱔⱅⱏ ⱄⱔ ⱄⰲⱁⰱⱁⰴⱐⱀⰻ ⰻ ⱃⰰⰲⱐⱀⰻ ⰲⱏ ⰴⱁⱄⱅⱁⰻⱀⱐⱄⱅⰲⱑ ⰻ ⰸⰰⰽⱁⱀⱑ}
    \end{block}

    \begin{OneImage}[htbp][\linewidth][0.25\textheight]{_assets/20260310_130957.png}[格拉哥里文手寫體（Cursive）]\end{OneImage}
\end{frame}

\begin{frame}{西里爾文}
    \textbf{西里爾文}（Cyrillic）的創制稍晚於格拉哥里文，都用於書寫古教會斯拉夫語（Old Church Slavic）。
    傳統上將西里爾文的創制也歸因於基里爾，但這一點存在爭議。
    相對比較明確的是，西里爾文更大程度地借用了希臘文，相對於格拉哥里文來說更加簡便，不斷擴張，最後成爲了斯拉夫世界的主流。

    \begin{OneImage}[htbp][\linewidth][0.5\textheight]{_assets/20260310_135648.png}[早期西里爾文]\end{OneImage}
\end{frame}

\begin{frame}{西里爾文}
    西里爾文在創制後經歷過多次正字法改革，其中最重要的是1708年彼得一世改革與1918年俄語正字法改革。

    \begin{block}{1708年彼得一世正字法改革}
        在中世紀，西里爾文演變出諸多地方變體。1708年，爲了推動世俗教育，規範世俗文本的西里爾文字體，彼得一世（\GK{Пётр Первый}）簡化了教會書寫，引入更簡便、現代的書寫形式，建立並推行了俄語民用體（\GK{Гражданский шрифт}）。這一套民用體，成爲了後世俄語書寫的主要來源。
    \end{block}
    
    \begin{block}{1918年俄語正字法改革}
        十月革命後，爲了簡化正字法，適應語言演變，蘇維埃政府推行了1918年俄語正字法改革（\GK{Реформа русской орфографии 1918 года}）。這次改革重點在於使得俄語正字法更符合俄語的實際發音，其中最具代表性的即是剔除了純粹形式性的字母\GK{ъ}。這次改革塑造了現代俄語西里爾文，沿用至今。
    \end{block}
\end{frame}

\begin{frame}{西里爾文}
    我們選取現代俄語來對西里爾文進行介紹。
    
    \small
    {\setlength{\columnsep}{0.6em}% ← 三欄之間距離：0.4em~0.8em 自己微調
    \setlength{\tabcolsep}{3pt}%  ← 表內欄距：2pt~4pt 自己微調

    \begin{columns}[T]
        \column{0.3\textwidth}
            \centering
            \begin{tabularx}{\linewidth}{@{}lllc@{}}
                \toprule
                \NiceHead{ 字母 & 名稱 & IPA & 對照 \\ }
                \GK{А а} & \GK{а} & /a/ & \GL{Ⰰ} \\
                \GK{Б б} & \GK{бэ} & /b/ & \GL{Ⰱ} \\
                \GK{В в} & \GK{вэ} & /v/ & \GL{Ⰲ} \\
                \GK{Г г} & \GK{гэ} & /g/ & \GL{Ⰳ} \\
                \GK{Д д} & \GK{дэ} & /d/ & \GL{Ⰴ} \\
                \GK{Е е} & \GK{е} & /je/ & \GL{Ⰵ} \\
                \GK{Ё ё} & \GK{ё} & /jo/ & \GL{Ⱖ} \\
                \GK{Ж ж} & \GK{жэ} & /ʐ/ & \GL{Ⰶ} \\
                \GK{З з} & \GK{зэ} & /z/ & \GL{Ⰸ} \\
                \GK{И и} & \GK{и} & /i/ & \GL{Ⰻ} \\
                \GK{Й й} & \GK{и краткое} & /j/ & \GL{Ⰺ} \\

                \bottomrule
            \end{tabularx}

        \column{0.25\textwidth}
            \centering
            \begin{tabularx}{\linewidth}{@{}lllc@{}}
                \toprule
                \NiceHead{ 字母 & 名稱 & IPA & 對照 \\ }
                \GK{К к} & \GK{ка} & /k/ & \GL{Ⰽ} \\
                \GK{Л л} & \GK{эль} & /l/ & \GL{Ⰾ} \\
                \GK{М м} & \GK{эм} & /m/ & \GL{Ⰿ} \\
                \GK{Н н} & \GK{эн} & /n/ & \GL{Ⱀ} \\
                \GK{О о} & \GK{о} & /o/ & \GL{Ⱁ} \\
                \GK{П п} & \GK{пэ} & /p/ & \GL{Ⱂ} \\
                \GK{Р р} & \GK{эр} & /r/ & \GL{Ⱃ} \\
                \GK{С с} & \GK{эс} & /s/ & \GL{Ⱄ} \\
                \GK{Т т} & \GK{тэ} & /t/ & \GL{Ⱅ} \\
                \GK{У у} & \GK{у} & /u/ & \GL{Ⱆ} \\
                \GK{Ф ф} & \GK{эф} & /f/ & \GL{Ⱇ} \\

                \bottomrule
            \end{tabularx}

        \column{0.4\textwidth}
            \centering
            \begin{tabularx}{\linewidth}{@{}lllc@{}}
                \toprule
                \NiceHead{ 字母 & 名稱 & IPA & 對照\\ }
                \GK{Х х} & \GK{ха} & /x/ & \GL{Ⱈ} \\
                \GK{Ц ц} & \GK{цэ} & /ts/ & \GL{Ⱌ} \\
                \GK{Ч ч} & \GK{че} & /tɕ/ & \GL{Ⱍ} \\
                \GK{Ш ш} & \GK{ша} & /ʂ/ & \GL{Ⱎ} \\
                \GK{Щ щ} & \GK{ща} & /ɕ/ & \GL{Ⱋ} \\
                \GK{Ъ ъ} & \GK{твёрдый знак} & -- & \GL{Ⱏ} \\
                \GK{Ы ы} & \GK{ы} & /ɨ/ & \GL{ⰟⰊ} \\
                \GK{Ь ь} & \GK{мягкий знак} & /ʲ/ & \GL{Ⱐ} \\
                \GK{Э э} & \GK{э} & /ɛ/ & \GL{Ⰵ} \\
                \GK{Ю ю} & \GK{ю} & /ju/ & \GL{Ⱓ} \\
                \GK{Я я} & \GK{я} & /ja/ & \GL{Ⱑ} \\
                \bottomrule
            \end{tabularx}
    \end{columns}
    }
\end{frame}

\begin{frame}{西里爾文文本案例}
    \begin{tabular}{ll}
        \GK{\textbf{Я вас любил} -- А. Пушкин} & \textbf{我曾經愛過你}~\small{普希金}\\
        \GK{Я вас любил: любовь ещё, быть может,} & 我曾經愛過你：愛情，也許\\
        \GK{В душе моей угасла не совсем;} & 在我的心裏還沒有完全熄滅；\\        
        \GK{Но пусть она вас больше не тревожит;} & 但願它不再使你憂傷；\\    
        \GK{Я не хочу печалить вас ничем.} & 我不願再用任何事使你悲愁。\\
        \\
        \GK{Я вас любил безмолвно, безнадежно,} & 我曾經默默地、絕望地愛過你，\\
        \GK{То робостью, то ревностью томим;} & 一會兒爲怯懦，一會兒爲嫉妒所苦；\\
        \GK{Я вас любил так искренно, так нежно,} & 我曾經那樣真誠、那樣溫柔地愛過你，\\
        \GK{Как дай вам Бог любимой быть другим.} & 但願上帝保佑你，也有人像我一樣愛你。\\
    \end{tabular}
\end{frame}

\end{document}
